% sections/applications.tex
% Included from main.tex via \input{sections/applications}

\section{Representative Applications}
\label{sec:apps}

CARLA-Air is validated through five representative workflows that collectively
exercise the platform's core capabilities across the four research directions
identified in Section~\ref{sec:introduction}: air-ground cooperation, embodied
navigation and vision-language action, multi-modal perception and dataset
construction, and reinforcement-learning-based policy training.
Table~\ref{tab:workflows_summary} provides a structured overview; detailed
descriptions follow in Sections~\ref{sec:apps:coop}--\ref{sec:apps:rl}.

\input{figures/tab_workflows_summary}

All workflows share a common dual-client architecture: both API clients execute
within the same Python process and operate on the same world state without
inter-process communication.
The following minimal pattern is common to every workflow:

{\footnotesize
\begin{verbatim}
ground = carla.Client('localhost', 2000)
aerial = airsim.MultirotorClient()
world  = ground.get_world()  # shared
aerial.enableApiControl(True)
\end{verbatim}
}

\input{figures/fig_workflow_arch}

%% ---------------------------------------------------------------
\subsection{W1: Air-Ground Cooperative Precision Landing}
\label{sec:apps:coop}

Autonomous precision landing of a UAV on a moving ground vehicle is a
representative and challenging scenario for air-ground cooperation, with
direct applications in drone-assisted logistics, mobile recharging, and
multi-agent coordination.
This workflow demonstrates real-time cross-domain coordination within
CARLA-Air's tick-synchronous control loop: the drone must continuously track
the vehicle's position, plan a descent trajectory, and execute a smooth
touchdown---all while the vehicle is in motion through urban traffic.

\paragraph{Setup.}
A ground vehicle follows an autopilot route through Town10HD at moderate
speed.
A drone starts at ${\approx}12$\,m altitude above the vehicle and is tasked
with landing on the vehicle's roof.
At each synchronous tick, the vehicle's 3D pose is queried via the ground API,
transformed into the aerial NED frame using the coordinate mapping from
Section~\ref{sec:arch:coords}, and used to compute a descent command issued
to the aerial flight controller.

\paragraph{Control architecture.}
The landing controller operates in three phases: \emph{approach}
(horizontal alignment with the vehicle), \emph{descent} (controlled altitude
reduction while maintaining horizontal tracking), and \emph{touchdown}
(final landing).
Let $\mathbf{q}_k \in \mathbb{R}^2$ denote the vehicle's horizontal position
at tick $k$ in the NED frame.
The drone's commanded horizontal target tracks the vehicle continuously:
%
\begin{equation}
\hat{\mathbf{q}}_k = \mathbf{q}_k + \mathbf{d},
\qquad
\hat{z}_k = z_k^{\mathrm{ref}},
\label{eq:landing_control}
\end{equation}
%
where $\mathbf{d}$ is the coordinate origin offset from
Eq.~\eqref{eq:coord_offset} and $z_k^{\mathrm{ref}}$ decreases according to
a smooth descent profile from the initial altitude to ground level.

\paragraph{Results.}
Fig.~\ref{fig:w1_coop} shows the complete landing sequence.
The drone descends from ${\approx}12$\,m to touchdown over ${\approx}20$\,s,
with horizontal convergence error decreasing from ${\approx}6$\,m to within
the $\pm 0.5$\,m tolerance band.
The 3D trajectory overview confirms that the UAV trajectory converges smoothly
onto the vehicle trajectory throughout the descent.
Table~\ref{tab:coop_results} summarizes the landing performance.

\input{figures/tab_coop_results}

\begin{figure*}[!htbp]
\centering
\includegraphics[width=\textwidth, height=0.38\textheight, keepaspectratio]{figures/landing.jpg}
\caption{%
  W1: Air-ground cooperative precision landing on a moving vehicle.
  \textbf{(a)}~Time-lapse sequence ($t_1$--$t_5$) showing the drone
  (red box) descending toward and landing on a moving ground vehicle
  through approach, descent, and touchdown phases.
  \textbf{(b)}~3D trajectory overview: the UAV trajectory (blue) converges
  onto the vehicle trajectory (orange), with the ground projection
  (dashed) showing horizontal alignment.
  \textbf{(c)}~Altitude profile over time, illustrating smooth descent from
  12\,m to ground level.
  \textbf{(d)}~Horizontal convergence error, decreasing from
  ${\approx}6$\,m to within the $\pm 0.5$\,m tolerance band.
  All data are recorded within CARLA-Air's synchronous tick loop.%
}
\vspace{-2mm}
\label{fig:w1_coop}
\end{figure*}

%% ---------------------------------------------------------------
\subsection{W2: Embodied Navigation and VLN/VLA Data Generation}
\label{sec:apps:vln}

Vision-language navigation (VLN) and vision-language action (VLA) are among the
most active research directions in embodied intelligence, requiring agents to
navigate and act in realistic environments grounded in natural language
instructions and visual observations.
A key bottleneck is the availability of large-scale, diverse training data that
pairs language instructions with rich visual observations from multiple
viewpoints.
CARLA-Air provides a natural data generation infrastructure for this purpose:
its photorealistic urban environments, socially-aware pedestrians, dynamic
traffic, and simultaneous aerial-ground sensing enable the construction of
VLN/VLA datasets with cross-view visual grounding that single-domain platforms
cannot provide.

\paragraph{Platform capabilities for VLN/VLA.}
CARLA-Air supports VLN/VLA data generation through several platform-level
features:
(i)~both aerial and ground agents can be equipped with egocentric RGB, depth,
and semantic segmentation cameras, providing paired bird's-eye and street-level
visual observations along any navigation trajectory;
(ii)~the ground API exposes waypoint-based route planning with lane-level
precision, which can serve as the basis for generating language-grounded
navigation instructions;
(iii)~the shared rendering pipeline guarantees that all visual observations are
captured under identical weather, lighting, and scene conditions at each tick;
and (iv)~the aerial overview provides a natural ``oracle view'' for generating
spatial referring expressions and verifying navigation progress.



\begin{figure*}[!htbp]
\centering
\includegraphics[width=\textwidth]{figures/vln.jpg}
\caption{%
  W2: Embodied navigation with aerial reasoning.
  A UAV autonomously tracks a pedestrian (red box, bottom) through an urban
  environment using bird's-eye visual observations.
  Each frame is annotated with the drone's chain-of-thought reasoning, illustrating how the agent interprets the scene,
  anticipates occlusions, and adjusts its flight path to maintain persistent
  visual contact with the target.
  All frames are rendered within CARLA-Air's shared simulation world under
  consistent lighting and physics.%
}
\label{fig:w2_vln}
\end{figure*}

%% ---------------------------------------------------------------
\subsection{W3: Synchronized Multi-Modal Dataset Collection}
\label{sec:apps:dataset}

Generating large-scale paired aerial-ground datasets is a bottleneck for
training and evaluating cooperative perception models.
Manual synchronization across separate simulator processes introduces alignment
errors that corrupt correspondence annotations.
CARLA-Air eliminates this bottleneck: because both sensor suites are driven by
the same tick counter, the resulting dataset records carry guaranteed per-tick
correspondence with zero interpolation overhead.

\paragraph{Setup.}
Eight ground sensors (RGB, semantic segmentation, depth, LiDAR, radar, GNSS,
IMU, collision) and four aerial sensors (RGB, depth, IMU, GPS) are attached
concurrently.
The simulation runs in synchronous mode at a fixed timestep; all 12 sensor
callbacks are registered before the first tick advance.
Ground traffic is populated with 30 autopilot vehicles and 10 pedestrians.

\paragraph{Dataset structure.}
Each record $\mathcal{R}_k$ at tick $k$ contains all 12 sensor observations
sharing a common tick index, plus vehicle and drone pose in the unified world
frame.
Records are serialized to disk in a flat per-tick directory structure with
metadata in JSON.
No timestamp interpolation is required: the shared tick index serves as the
alignment key.

\paragraph{Results.}
Over 1\,000 ticks, the workflow produces 1\,000 fully synchronized 12-stream
records at a mean collection rate of ${\approx}17$\,Hz.
The maximum observed tick-alignment deviation across all 12 streams is one tick
(occurring transiently under disk-write saturation).
Table~\ref{tab:dataset_results} summarizes collection performance.

\input{figures/tab_dataset_results}


% Replace the W3 figure placeholder in sections/applications.tex with:

\begin{figure*}[!htbp]
\centering
\includegraphics[width=\textwidth]{figures/datasets.jpg}
\caption{%
  W3: Synchronized multi-modal dataset collection at a single simulation tick.
  \textbf{Top row} (vehicle perspective): RGB, semantic segmentation, depth,
  LiDAR bird's-eye view, surface normals, and instance segmentation.
  \textbf{Bottom row} (drone perspective): the same six modalities captured
  simultaneously from the aerial viewpoint.
  All 12 sensor streams share an identical tick index and are rendered under
  the same lighting and weather conditions through CARLA-Air's shared
  rendering pipeline, guaranteeing per-tick spatial-temporal correspondence
  without interpolation.%
}
\label{fig:w3_dataset}
\end{figure*}

%% ---------------------------------------------------------------
\subsection{W4: Air-Ground Cross-View Perception}
\label{sec:apps:perception}

Cross-view perception---fusing aerial bird's-eye observations with ground-level
sensing---is an emerging research direction for cooperative autonomous driving,
urban scene understanding, and 3D reconstruction.
This workflow demonstrates CARLA-Air's ability to provide spatially and
temporally co-registered multi-modal sensor streams from both aerial and ground
viewpoints within a single shared environment.

\paragraph{Setup.}
A drone equipped with a depth camera hovers above a road segment while a ground
ego vehicle equipped with a semantic segmentation camera traverses the same
segment.
Both sensors are queried at each synchronous tick.

\paragraph{Sensor co-registration.}
Co-registration leverages the coordinate mapping from
Section~\ref{sec:arch:coords}.
The origin offset $\mathbf{d} \in \mathbb{R}^3$ between the aerial NED frame
and the ground world frame is computed once at initialization:
%
\begin{equation}
\mathbf{d} = T\!\left(\mathbf{p}_{\mathrm{spawn}}^{\mathrm{world}}\right)
  - \mathbf{p}_{\mathrm{spawn}}^{\mathrm{NED}},
\label{eq:coord_offset}
\end{equation}
%
where $T(\cdot)$ applies the scale conversion and axis remapping.
When the drone spawns at the world origin, $\mathbf{d} = \mathbf{0}$.

\paragraph{Weather consistency.}
To verify rendering consistency across both sensor layers, the workflow iterates
through all 14 official weather presets.
For each preset involving significant illumination change, the mean pixel
intensity of the aerial RGB frame shifts by more than 5\% relative to the
previous preset, confirming single-pass weather propagation.
Because both sensors share the same rendering pipeline, temporal alignment is
guaranteed:
%
\begin{equation}
\epsilon_k = \bigl|t_k^{\mathrm{gnd}} - t_k^{\mathrm{air}}\bigr| = 0,
\label{eq:sensor_alignment}
\end{equation}
%
a property that bridge-based architectures cannot provide.

\paragraph{Results.}
Over 500 ticks, the workflow produces 500 co-registered aerial-depth /
ground-segmentation pairs at ${\approx}18$\,Hz with zero RPC errors.
All 14 weather presets pass the illumination consistency assertion.
Table~\ref{tab:perception_results} summarizes the measured outcomes.

\input{figures/tab_perception_results}



\begin{figure*}[!t]
\centering
\includegraphics[width=\textwidth, height=0.32\textheight, keepaspectratio]{figures/w4.jpg}
\caption{%
  W4: Air-ground cross-view perception across diverse environments and weather
  conditions. Each row shows an aerial RGB view from the drone across six
  representative CARLA maps (Town\,01--05 and Town\,10HD); each column
  corresponds to a different weather preset (Clear Noon, Cloudy Noon, Dense
  Fog, Hard Rain, Night, Soft Rain, and Sunset).%
}
\label{fig:w4_perception}
\end{figure*}

%% ---------------------------------------------------------------
\subsection{W5: Reinforcement Learning Training Environment}
\label{sec:apps:rl}

Reinforcement learning in air-ground cooperative settings requires a simulation
environment that provides closed-loop interaction, consistent state
observations across heterogeneous agents, and stable long-horizon episode
execution without memory leaks or synchronization drift.
CARLA-Air's single-process architecture naturally satisfies these requirements,
and the stability results from Section~\ref{sec:perf:vram} (zero crashes and
zero memory accumulation over 357 reset cycles) directly validate its
suitability as an RL training environment.

\paragraph{Platform capabilities for RL.}
CARLA-Air supports RL training workflows through several platform-level
features:
(i)~synchronous stepping mode provides deterministic state transitions
compatible with standard Gym-style training loops;
(ii)~both aerial and ground agents expose programmatic control interfaces (velocity
commands, waypoint targets, autopilot toggles) that can serve as action spaces;
(iii)~the full sensor suite (RGB, depth, segmentation, LiDAR, IMU, GPS)
provides rich observation spaces for both state-based and vision-based policies;
(iv)~episode resets via actor spawn/destroy are validated for stability across
hundreds of consecutive cycles (Section~\ref{sec:perf:vram});
and (v)~the shared world state ensures that reward signals computed from
cross-domain interactions (e.g., aerial-ground relative positioning) are
physically consistent.

\paragraph{Example: cooperative positioning.}
As a representative RL scenario, consider a drone learning to maintain an
optimal aerial observation position relative to a moving ground vehicle under
varying traffic conditions.
The observation space comprises the drone's pose, the vehicle's pose, and
surrounding traffic state; the action space is a 3D velocity command; the reward
encodes lateral tracking error and altitude maintenance.
This scenario exercises the full air-ground control loop within CARLA-Air's
synchronous tick, and can be implemented using standard RL libraries (e.g.,
Stable-Baselines3, RLlib) with minimal wrapper code.


\begin{figure*}[!htbp]
\centering
\includegraphics[width=\textwidth, height=0.35\textheight, keepaspectratio]{figures/RL.jpg}
\caption{%
  W5: Reinforcement learning training environment.
  \textbf{Top:} a drone learns to maintain an optimal aerial observation
  position above a moving ground vehicle within CARLA-Air's urban traffic
  environment.
  \textbf{Bottom:} the closed-loop RL pipeline. At each synchronous tick,
  CARLA-Air provides an observation space (drone and vehicle poses, relative
  distance, surrounding traffic state) to the policy network, which outputs
  3D velocity commands as actions. The reward signal encodes tracking
  accuracy, altitude maintenance, and collision avoidance.%
}
\label{fig:w5_rl}
\end{figure*}